\chapter[Descente des morphismes \'etales]{Descente des morphismes \'etales.\\ Application au groupe fondamental}
\label{IX}

\section{Rappels sur les morphismes \'etales}
\label{IX.1}

Nous allons ici passer en revue les propri\'et\'es des morphismes
\'etales (d\'evelopp\'es dans~\ref{I}) qui vont nous servir,
en profitant de cette occasion pour \'eliminer de la th\'eorie les
hypoth\`eses noeth\'eriennes superflues. Le lecteur notera que
m\^eme si on ne s'int\'eresse qu'aux sch\'emas noeth\'eriens, la
technique de descente conduit \`a introduire des sch\'emas non
noeth\'eriens (tels que~$\Spec(\widehat{A}\otimes_{A}\widehat{A})$,
o\`u~$A$ est un anneau local noeth\'erien), et pour pouvoir
appliquer le langage des cat\'egories fibr\'ees, il importe de
d\'efinir les notions telles que morphisme \'etale etc., sans y
introduire de restriction noeth\'erienne. Le lecteur qui
r\'epugnerait \`a v\'erifier ou \`a admettre que les
\'enonc\'es ci-dessous sont vrais sans hypoth\`eses
noeth\'eriennes, pourra se contenter de les admettre sous les
hypoth\`eses noeth\'eriennes de~\ref{I}, \`a condition
d'introduire ces m\^emes hypoth\`eses noeth\'eriennes dans les
\'enonc\'es des num\'eros suivants, et d'utiliser la
d\'efinition~\ref{IX.1.1} ci-dessous pour les sch\'emas non
noeth\'eriens qui s'introduisent dans les raisonnements.

\begin{definition}
\label{IX.1.1}
Soient~$f\colon X\to S$ un morphisme de pr\'esch\'emas, et~$x$ un
point de~$X$. On dit que~$f$ est \emph{étale en}~$x$,
\index{etale (morphisme, algebre)@\'etale (morphisme, alg\`ebre)|hyperpage}ou que~$X$ est \'etale sur~$S$ en~$x$, s'il existe un voisinage
ouvert affine~$U$ de~$s=f(x)$, un voisinage ouvert affine~$V$ de~$x$
au-dessus de~$U$, un sch\'ema \emph{noethérien} affine~$U_{0}$,
un~$U_{0}$-sch\'ema \emph{étale}~\eqref{I} et affine~$V_{0}$, un
morphisme~$U\to U_{0}$, et un~$U$-isomorphisme
$$
V\isomto V_{0}\times_{U_{0}}U\quoi.
$$
\end{definition}

On
notera que lorsque~$S$ est localement noeth\'erien, cette
terminologie co\"incide avec celle de \loccit On dira de m\^eme
que~$f$ \emph{est étale}, ou que~$X$ \emph{est étale sur}~$S$,
si~$f$ est \'etale en tout point~$x$ de~$X$. Avec ces
d\'efinitions, les propositions ci-dessous se ram\`enent sans
difficult\'e au cas noeth\'erien, o\`u elles sont
d\'emontr\'ees dans~I~\No \ref{I.4},
\ref{I.5},~\ref{I.7}. Pour des d\'etails, le lecteur pourra
consulter EGA~IV\footnote{De fa\c con pr\'ecise, EGA~IV 17,~18.}.

\begin{remarques}
\label{IX.1.2}
Si $f$ est \'etale en $x$, alors $f$ est \og \emph{de présentation
finie en}~$x$\fg (VIII~\ref{VIII.3.5}), l'anneau local de~$x$ dans la
fibre $f^{-1}(s)$ est une \emph{extension finie séparable} de
$\kres(s)$, enfin $f$ est \emph{plat en}~$x$. On peut montrer que la
r\'eciproque est vraie, donc que la d\'efinition~\ref{IX.1.1} est
la m\^eme que dans le cas o\`u $S$ est localement noeth\'erien,
sauf qu'il faut remplacer la condition \og de type fini en~$x$\fg par
\og de pr\'esentation finie en~$x$\fg. Comme ce r\'esultat est de
d\'emonstration d\'elicate, nous n'avons pas voulu ici donner
cette d\'efinition de la notion de morphisme \'etale, qui ne se
pr\^ete pas directement \`a la d\'emonstration des
propri\'et\'es qui vont suivre.
\end{remarques}

Notons d'abord qu'on a trivialement:
\begin{proposition}
\label{IX.1.3}
Si $f\colon X\to S$ est \'etale, alors tout morphisme $f'\colon
X'\to S'$ qui s'en d\'eduit par changement de base $S'\to S$ est
\'egalement \'etale.
\end{proposition}

On peut donc dire que les morphismes \'etales forment une
\emph{sous-catégorie fibrée} de la cat\'egorie des
fl\`eches dans~$\Sch$
(\cf VI~\ref{VI.11}~a)). L'objet du pr\'esent
expos\'e est l'\'etude des propri\'et\'es d'exactitude de
cette cat\'egorie fibr\'ee sur~$\Sch$.
\begin{proposition}
\label{IX.1.4}
Soit $f\colon X\to S$ un morphisme de pr\'esch\'emas. Pour que ce
soit une immersion ouverte, il faut et il suffit qu'il soit \'etale
et radiciel.
\end{proposition}

\Cf I~\ref{I.5.1}. On en conclut que si $X$ est \'etale sur~$S$, toute section de~$X$ sur~$S$ est une immersion ouverte, donc,
utilisant encore~\ref{IX.1.4}, on trouve:
\begin{corollaire}
\label{IX.1.5}
Soit $X$ un $S$-pr\'esch\'ema \'etale. Alors il y a une
correspondance biunivoque entre l'ensemble des sections de~$X$
sur~$S$, et l'ensemble des parties
ouvertes $\Gamma$ de~$X$ telles que le morphisme $\Gamma \to S$ induit
par le morphisme structural soit \emph{radiciel} et \emph{surjectif}.
\end{corollaire}

Si d'ailleurs $X$ est s\'epar\'e sur~$S$, $\Gamma$ sera une partie
de~$X$ \`a la fois ouverte et ferm\'ee, mais peu importe. --- Faisant un changement de base \'evident, on peut mettre~\ref{IX.1.5}
sous la forme en apparence plus g\'en\'erale:
\begin{corollaire}
\label{IX.1.6}
Soient $X$ et $Y$ deux $S$-pr\'esch\'emas, $Y$ \'etant \'etale
sur~$S$. Alors l'application $f \mto \Gamma_f$ qui associe \`a
tout $S$-morphisme $f$ de~$X$ dans $Y$ la partie de~$X\times_S Y$
sous-jacente au graphe de~$f$, est une bijection de~$\Hom_S(X,Y)$ sur
l'ensemble des parties ouvertes de~$\Gamma$ de~$X\times_S Y$ telles
que le morphisme $\Gamma\to X$ induit par $\pr_1$ soit
\emph{radiciel} et \emph{surjectif}.
\end{corollaire}

\begin{proposition}
\label{IX.1.7}
Soit $S_0$ le sous-pr\'esch\'ema de~$S$ d\'efini par un
Nil-id\'eal quasi-coh\'erent, \ie tel que $S_0$ ait m\^eme
ensemble sous-jacent que~$S$. Alors le foncteur $X \mto X\times_S
S_0$ de la cat\'egorie des pr\'esch\'emas \'etales sur~$S$
dans la cat\'egorie des pr\'esch\'emas \'etales sur~$S_0$, est
une \'equivalence de cat\'egories.
\end{proposition}

Le fait que ce foncteur soit pleinement fid\`ele est une
cons\'equence imm\'ediate de~\ref{IX.1.6}. Le fait qu'il soit
essentiellement surjectif est contenu dans~I~\ref{I.8.3}. On notera
que dans l'\'equivalence pr\'ec\'edente,
$X$ est de type fini \ie quasi-fini sur~$S$, (\resp fini \ie un rev\^etement \'etale
de~$S$), si et seulement si $X_0$ satisfait \`a la condition
analogue sur~$S_0$; m\^eme remarque pour la condition de
s\'eparation. Ces faits sont imm\'ediats, et aussi contenus
dans~\ref{IX.2.4} plus bas.

\begin{corollaire}
\label{IX.1.8}
Soit $A$ un anneau local noeth\'erien complet de corps r\'esiduel
$k$. Alors le foncteur $B \mto B\otimes_A k$ est une
\'equivalence de la cat\'egorie des alg\`ebres finies et
\'etales sur~$A$, avec la cat\'egorie des alg\`ebres finies et
\'etales sur~$k$, (\ie compos\'ees d'un nombre fini d'extensions
finies s\'eparables de~$k$).
\end{corollaire}

\begin{proposition}
\label{IX.1.9}
Pour que $X$ soit un \emph{revêtement étale} de~$S$ \ie fini
et \'etale sur~$S$, il faut et il suffit que $X$ soit $S$-isomorphe
au spectre d'une Alg\`ebre $\cal{A}$ sur~$S$, qui soit un Module
localement libre de type fini, et telle que pour tout
$s \in S$, $\cal{A}_s\otimes_{\cal{O}_s}\kres(s)$ soit une alg\`ebre
s\'eparable sur~$\kres(s)$, donc en l'occurrence compos\'ee directe
d'extensions finies s\'eparables de~$\kres(s)$.
\end{proposition}

Enfin le r\'esultat suivant est de nature moins \'el\'ementaire,
\'etant la conjonction de~I~\ref{I.8.4} et du
\emph{théorème d'existence de faisceaux en géométrie
algébrique}, (EGA~III~5; \cf aussi \cite{IX.1} th.\, 3).

\begin{theoreme}
\label{IX.1.10}
Soient $S$ le spectre d'un anneau local noeth\'erien complet, $X$ un
$S$-sch\'ema propre, $X_0$ la fibre de~$X$ au point ferm\'e de
$S$, (de sorte que $X_0$ est un sous-sch\'ema ferm\'e de
$X$). Alors le foncteur restriction $X' \mto X'\times_X X_0$ est
une \'equivalence de la cat\'egorie des rev\^etements \'etales
de~$X$ avec la cat\'egorie des rev\^etements \'etales de~$X_0$.
\end{theoreme}

\section{Morphismes submersifs et universellement submersifs}
\label{IX.2}

\begin{definition}
\label{IX.2.1}
Un morphisme $g\colon S' \to S$ de pr\'esch\'emas est dit
\emph{submersif}
\index{submersif (morphisme)|hyperpage}s'il est surjectif, et fait de~$S$ un espace topologique quotient
de~$S'$ (\ie une partie $U$ de~$S$ telle que $f^{-1}(U)$ soit
ouverte, est ouverte). On dit que $f$ est \emph{universellement
submersif}
\index{universellement submersif (morphisme)|hyperpage}si pour tout morphisme $T \to S$, le morphisme $f'\colon T'=S'\times_S
T \to T$ d\'eduit de~$f$ par changement de base est submersif.
\end{definition}

Il est imm\'ediat que le compos\'e de deux morphismes submersifs
(\resp universellement submersifs) est submersif
(\resp universellement submersif), et qu'un changement de base dans un
morphisme universellement submersif donne un morphisme universellement
submersif (vu qu'on fait ce qu'il faut pour cela). Si $fg$ est
submersif (\resp universellement submersif), $f$ l'est.
\begin{exemples}
\label{IX.2.2}
a) Un morphisme surjectif qui est ouvert, ou ferm\'e, est submersif,
donc un morphisme surjectif universellement ferm\'e ou
universellement ouvert est universellement submersif. Par exemple
\emph{un morphisme propre surjectif est universellement
submersif}. D'autre part \emph{un morphisme fidèlement plat et
quasi-compact est universellement submersif} (VIII~\ref{VIII.4.3}). Ce
seront les deux cas les plus importants pour nous.
\end{exemples}

On
peut appliquer \`a un morphisme submersif ou universellement
submersif $g\colon S' \to S$ les raisonnements de~VIII~\ref{VIII.4.3},
on trouve en particulier:
\begin{proposition}
\label{IX.2.3}
Supposons $g\colon S' \to S$ submersif. Alors le diagramme suivant
d'applications est exact:
$$
\xymatrix@C=.5cm{{\mathrm{Ouv} (S)}\ar[r]&{\mathrm{Ouv} (S')}\ar@<2pt>[r]\ar@<-2pt>[r]&{\mathrm{Ouv}(S''),}}
$$
o\`u $S'' = S'\times_S S'$, et o\`u $\Ouv(X)$ d\'esigne
l'ensemble des parties ouvertes du pr\'esch\'ema~$X$.
\end{proposition}

\begin{proposition}
\label{IX.2.4}
Soient $g\colon S' \to S$ un morphisme universellement submersif,
$f\colon X \to Y$ un $S$-morphisme, et $f'\colon X' \to Y'$ le
$S'$-morphisme qui s'en d\'eduit par changement de base. Pour que
$f$ soit ouverte (\resp ferm\'ee), il suffit que $f'$ le soit. Pour
que $f$ soit universellement ouverte, \resp universellement
ferm\'ee, \resp s\'epar\'ee, il faut et il suffit que $f'$ le
soit. Si de plus $g$ est quasi-compact, et $f$ localement de type
fini, pour que $f$ soit propre, il faut et il suffit que $f'$ le soit.
\end{proposition}

Pour ce dernier point, on note que si $f'$ est propre donc
quasi-compact alors $f$ est quasi-compact (VIII~\ref{VIII.3.3}) donc
de type fini puisqu'il est localement de type fini. D'autre part il
est s\'epar\'e et universellement ferm\'e d'apr\`es ce qui
pr\'ec\`ede, donc il est propre.
\begin{proposition}
\label{IX.2.5}
Soit $S'$ un pr\'esch\'ema de type fini sur le spectre $S$ d'un
anneau local noeth\'erien complet, supposons que la fibre du point
ferm\'e $s$ de~$S$ soit finie, donc les anneaux locaux dans $S'$ des
points $s'$ de cette fibre sont finis sur~$A=\cal{O}_s$. Soit~$S''$ le
sch\'ema somme des spectres des $\cal{O}_{S',s'}$ en question,
consid\'er\'e comme $S$-sch\'ema fini. Pour que $g\colon S' \to
S$ soit universellement submersif, il faut et il suffit que le
morphisme structural $S''\to S'$ soit surjectif.
\end{proposition}

Comme il y a un $S$-morphisme naturel $S'' \to S'$, et qu'un morphisme
fini surjectif est universellement submersif d'apr\`es~\ref{IX.2.2},
la condition \'enonc\'ee est suffisante. Montrons donc que si
$S''\to S'$ n'est pas surjectif, alors $g$ n'est pas universellement
submersif. En effet, soit $t$ un point de~$S$ qui n'est pas dans
l'image de~$S''$; il existe alors un $S$-sch\'ema $T$, spectre d'un
anneau de valuation discr\`ete, dont l'image dans $S$ est
$\{s,t\}$. Notons que l'image de~$S''$ dans $S'$ est ouverte, car le
morphisme $S'' \to S'$ est un isomorphisme local, et d'autre part
cette image contient $S'_s$, et ne rencontre pas $S'_t$. Il s'ensuit
que l'image inverse de cette derni\`ere dans $T'=S'\times_S T$ est
\emph{ouverte}, et identique \`a l'image inverse du point ferm\'e
de~$T$. Cela montre que $T'\to T$ n'est pas submersif, donc $S'\to S$
n'est pas universellement submersif.
\begin{remarque}
\label{IX.2.6}
Utilisant le crit\`ere IV~\ref{IV.6.3} pour qu'une partie
constructible d'un espace noeth\'erien soit ouverte, on trouve
facilement le crit\`ere valuatif suivant pour qu'un morphisme
$g\colon S'\to S$ \emph{de type fini}, avec $S$ localement
noeth\'erien, soit universellement submersif: il faut et il suffit
que pour tout $S$-sch\'ema $T$, spectre d'un anneau de valuation
discr\`ete, posant $T'=S'\times_S T$, l'image inverse dans $T'$ du
point ferm\'e de~$T$ soit non ouverte.
\end{remarque}

\section{Descente de morphismes de pr\'esch\'emas \'etales}
\label{IX.3}

\begin{proposition}
\label{IX.3.1}
Soient $g\colon S'\to S$ un morphisme \emph{surjectif} de
pr\'esch\'emas, $X$ et $Y$ deux pr\'esch\'emas sur~$S$, $X'$,
$Y'$ leurs images inverses sur~$S'$. Si $Y$ est non ramifi\'e sur~$S$, alors l'application canonique
$$
\Hom_S(X,Y) \to \Hom_{S'}(X',Y')
$$
est injective.
\end{proposition}

En effet, en vertu de~\ref{IX.1.6}, un $S$-morphisme $f\colon X\to Y$
est connu quand on conna\^it l'ensemble sous-jacent \`a son
graphe $\Gamma$, qui est une partie de~$Z=X\times_SY$. Comme
$$
Z'=Z\times_SS'=X'\times_{S'}Y' \to Z
$$
est surjectif, (puisque $S'\to S$ l'est), cette partie $\Gamma$ est
connue quand on conna\^it son image inverse dans $X'\times_{S'}Y'$,
qui n'est autre que l'ensemble sous-jacent au graphe de~$f'$. D'o\`u
la conclusion.

Une partie~$\Gamma$ de~$Z$ est le graphe d'un $S$-morphisme $f\colon
X\to Y$ si et seulement si elle est ouverte, et si le morphisme induit
par~$\pr_1$ de~$\Gamma$ dans~$Y$ est radiciel
et surjectif \cf \ref{IX.1}. Lorsque la premi\`ere
propri\'et\'e est v\'erifi\'ee, la deuxi\`eme l'est si et
seulement si l'image inverse~$\Gamma'$ de~$\Gamma$ dans~$Z'$ satisfait
la m\^eme condition VIII~\ref{VIII.3.1}. Si on sait enfin que
$Z'\to Z$ est submersif, ce qui sera le cas en particulier si $S'\to
S$ est universellement submersif, alors~$\Gamma$ est ouvert si et
seulement si~$\Gamma'$ l'est. Ainsi, l'ensemble $\Hom_S(X,Y)$ est
alors en correspondance biunivoque avec l'ensemble des parties
ouvertes~$\Gamma'$ de~$Z'$ telles que le morphisme projection
$\pr_1\colon Z'\to X'$ soit radiciel et surjectif, (\ie correspondant \`a un $S'$-morphisme $f'\colon X'\to Y'$), et qui
sont satur\'ees pour la relation d'\'equivalence d\'efinie par
$Z'\to Z$, \ie dont les deux images inverses dans $Z''=Z'\times_Z
Z'=Z\times_S S''$ (o\`u $S''=S'\times_S S'$), par l'une et l'autre
projection, sont \'egales. Or ces derni\`eres sont les graphes
des deux $S''$-morphismes $X''\to Y''$ d\'eduits de~$f'$ par
changement de base, par l'une et l'autre projection $S''\to S'$. On a
ainsi obtenu:

\begin{proposition}
\label{IX.3.2}
Soient $g\colon S'\to S$ un morphisme \emph{universellement submersif}
de pr\'esch\'emas, $S''=S'\times_S S'$, $X$ et~$Y$ deux
$S$-pr\'esch\'emas, $X'$ et~$Y'$ leurs images inverses sur~$S'$,
et $X''$, $Y''$ leurs images inverses sur~$S''$. Si~$Y$ est \'etale
sur~$S$ le diagramme canonique suivant d'applications est exact:
$$
\xymatrix@C=.5cm{{\mathrm{Hom}_S(X,Y)}\ar[r]&{\mathrm{Hom}_{S'}(X',Y')}\ar@<2pt>[r]\ar@<-2pt>[r]&{{\mathrm{Hom}_{S''}(X'',Y'')}.}}
$$
\end{proposition}

Prenant~$X$ et~$Y$ \'etales sur~$S$, on trouve l'\'enonc\'e
suivant, qui d'ailleurs redonne~\ref{IX.3.2} (m\^eme en se
restreignant \`a $X=S$, auquel cas on peut en effet toujours se
ramener dans~\ref{IX.3.2}, par le changement de base $X\to S$);

\begin{corollaire}
\label{IX.3.3}
Un morphisme universellement submersif de pr\'esch\'emas est un
morphisme de descente pour la cat\'egorie fibr\'ee des
pr\'esch\'emas \'etales sur d'autres.
\end{corollaire}

J'ignore d'ailleurs si c'est n\'ecessairement un morphisme de
descente \emph{effective} pour la cat\'egorie fibr\'ee en
question, m\^eme en faisant de plus l'hypoth\`ese que~$S$ est
noeth\'erien et~$g$ de type fini, et en se bornant aux
rev\^etements \'etales. Nous donnerons n\'eanmoins au
num\'ero suivant des crit\`eres utiles d'effectivit\'e.

\begin{corollaire}
\label{IX.3.4}
Soit
$g\colon S'\to S$ un morphisme universellement submersif, dont
les fibres $g^{-1}(s)$ sont \og g\'eom\'etriquement connexes\fg,
\index{geometriquement connexe@g\'eom\'etriquement connexe|hyperpage}\ie pour toute extension $K/\kres(s)$, $g^{-1}(s)\otimes_{\kres(s)} K$ est
connexe. Alors~$S'$ est connexe si~$S$ l'est. Le foncteur de la
cat\'egorie des pr\'esch\'emas \'etales sur~$S$ dans la
cat\'egorie
des pr\'esch\'emas \'etales sur~$S'$ d\'efini par~$g$ est
pleinement fid\`ele.
\end{corollaire}

Une partie de~$S'$ qui est \`a la fois ouverte et ferm\'ee est
satur\'ee pour la relation d'\'equivalence ensembliste d\'efinie
par~$g$, puisque les fibres sont connexes, donc est l'image inverse
d'une partie de~$S$, qui est n\'ecessairement ouverte et ferm\'ee
puisque~$g$ est submersif. Si donc~$S$ est connexe, $S'$ l'est. Cela
implique aussi le r\'esultat suivant: le compos\'e~$fg$ de deux
morphismes \`a fibres universellement connexes,
\index{universellement connexe: synonyme de g\'eom\'etriquement connexe|hyperpage}$f$ \'etant universellement submersif, est \`a fibres
universellement connexes; si~$S_1'$ et~$S_2'$ sur~$S$ ont des fibres
universellement connexes, il en est de m\^eme de~$S_1'\times_S
S_2'$. En particulier, sous les conditions de~\ref{IX.3.4}, $S''$ a
des fibres universellement connexes sur~$S$. Soient alors~$X$ et~$Y$
\'etales sur~$S$, et soit~$u'$ un $S'$-morphisme de~$X'$ dans~$Y'$,
prouvons qu'il est compatible avec les donn\'ees de descente (ce qui
entra\^ine la conclusion voulue gr\^ace \`a~\ref{IX.3.3}). Or
soient~$u_1''$ et~$u_2''$ les deux $S''$-morphismes $X''\to Y''$
d\'eduits de~$u'$. Le sous-pr\'esch\'ema de~$S''$ des
co\"incidences de~$u_1''$ et~$u_2''$ est un sous-pr\'esch\'ema
ouvert induit, ferm\'e fibre par fibre, comme image inverse du
pr\'esch\'ema diagonal de~$Y''$ sur~$S''$\footnote{noter que les
fibres de~$S'$ sur~$S$ sont s\'epar\'ees!}. C'est donc l'image
inverse d'une partie de~$S$. Comme elle contient la diagonale
dans~$S''$, elle est identique \`a~$S''$, d'o\`u $u_1''=u_2''$
cqfd.

\section{Descente de pr\'esch\'emas \'etales: crit\`eres
d'effectivit\'e}
\label{IX.4}

\enlargethispage{.7cm}\begin{proposition}
\label{IX.4.1}
Soit $g\colon S'\to S$ un morphisme fid\`element plat et
quasi-compact. Alors~$g$ est un morphisme de descente effective pour
la cat\'egorie fibr\'ee des pr\'esch\'emas \'etales,
s\'epar\'es et de type fini sur d'autres.
\end{proposition}

C'est en effet un morphisme de descente pour la cat\'egorie
fibr\'ee en question, en vertu de~\ref{IX.3.3} ou de
VIII~\ref{VIII.5.2} au choix. Reste \`a montrer que si~$X'$ est
\'etale, s\'epar\'e et de type fini sur~$S'$, et muni d'une
donn\'ee de descente relativement \`a $g\colon S'\to S$, cette
derni\`ere est effective dans la cat\'egorie fibr\'ee en
question. Or on voit facilement que si~$X$ est un pr\'esch\'ema
sur~$S$, alors
il est \'etale sur~$S$ si et seulement si il est \'etale sur~$S'$
(en vertu de la d\'efinition~\ref{IX.1.1} et de \loccit \ref{VIII.3.6}). Donc il est \'etale, s\'epar\'e et de
type fini sur~$S$ si et seulement si~$X'$ l'est sur~$S'$, \cf par
exemple~\ref{IX.2.4}. Donc il suffit de s'assurer de
l'effectivit\'e de la donn\'ee de descente sur~$X$ pour la
cat\'egorie fibr\'ee des fl\`eches de~$\Sch$. Or~$X'$ est
quasi-affine sur~$S'$ en vertu de VIII~\ref{VIII.6.2}
et~\ref{VIII.6.6}. On peut alors conclure en utilisant
VIII~\ref{VIII.7.9}. Le lecteur notera d'ailleurs que la
d\'emonstration demande moins si on se borne aux pr\'esch\'emas
\'etales et \emph{finis} sur d'autres, car on peut alors invoquer
directement VIII~\ref{VIII.2.1}.

\begin{corollaire}
\label{IX.4.2}
Soient $g\colon S'\to S$ un morphisme universellement submersif, $X'$
un
$S'$-pr\'esch\'ema
\'etale s\'epar\'e et de type fini, muni d'une donn\'ee de
descente relativement \`a~$g$, $S_1\to S$ un morphisme
fid\`element plat et quasi-compact, $S_1'$ et~$X_1'$ d\'eduits
de~$S'$ et~$X'$ par le changement de base, de sorte que $S_1'\to S_1$
est universellement submersif, $X_1'$ est \'etale s\'epar\'e et
de type fini sur~$S_1'$, et muni d'une donn\'ee de descente
relativement \`a $g_1\colon S_1'\to S_1$. Pour que la donn\'ee de
descente sur~$X'$ soit effective, il faut et il suffit que la
donn\'ee de descente sur~$X_1'$ le soit.
\end{corollaire}

Cela r\'esulte de la th\'eorie de la descente dans les
cat\'egories \cite{IX.D}, compte tenu de~\ref{IX.4.1} et~\ref{IX.3.3}.

On prouve de fa\c con analogue:

\begin{corollaire}
\label{IX.4.3}
Soient $g\colon S'\to S$ un morphisme universellement submersif, $X'$
un~$S'$-pr\'esch\'ema \'etale muni d'une donn\'ee de descente
relativement \`a~$g$, $(S_i)$ un recouvrement de~$S$ par des
ouverts. Pour que la donn\'ee de descente soit effective, il faut
et il suffit que pour tout~$i$, la donn\'ee de descente
correspondante sur~$X_i'=X\times_S S_i$, relativement au morphisme
$g_i\colon S_i'=S'\times_S S_i\to S_i$, le soit.
\end{corollaire}

Ce dernier r\'esultat conduit \`a d\'egager un crit\`ere
d'effectivit\'e local:

\enlargethispage{.3cm}
\begin{proposition}
\label{IX.4.4}
Soit $g\colon S'\to S$ un morphisme de pr\'esentation finie
\textup{VIII}~\ref{VIII.3.6}
et universellement submersif, $X'$ un pr\'esch\'ema \'etale et
de pr\'esentation finie sur~$S'$, muni d'une donn\'ee de descente
relativement \`a~$g$, enfin~$a$ un point de~$S$. Pour qu'il existe
un voisinage ouvert~$U$ de~$a$, tel que la donn\'ee de descente
correspondante sur~$X_U'=X'\times_S U$ relativement au morphisme
$$g_U\colon S_U'=S'\times_S U\to S_U=U$$ soit effective, il faut et il
suffit que la donn\'ee de descente correspondante sur $X_a'=X'\times_S\Spec(\cal{O}_a)$, relativement au morphisme
$$g_a\colon S_a'=S'\times_S\Spec(\cal{O}_a)\to S_a=\Spec(\cal{O}_a),$$
soit effective.
\end{proposition}

La n\'ecessit\'e \'etant triviale, montrons la suffisance. On
dispose donc d'un pr\'esch\'ema \'etale de type fini~$X_a$
sur~$S_a$, et d'un isomorphisme
\begin{equation*}
\label{eq:IX.4.*}
\tag{$*$} {X_a'\isomto X_a\times_{S_a}S_a'}
\end{equation*}
compatible avec les donn\'ees de descente. Conform\'ement \`a
un sorite g\'en\'eral facile sur les pr\'esch\'emas
d\'efinis sur une limite inductive d'anneaux (ici les anneaux~$A_f$,
o\`u~$A$ est l'anneau d'un voisinage ouvert affine de~$a$, et
o\`u~$f$ parcourt les \'el\'ements de~$A$ qui ne sont pas dans
l'id\'eal premier correspondant \`a~$a$), on peut trouver un
voisinage ouvert~$U$ de~$a$, un pr\'esch\'ema \'etale de type
fini~$X_U$ sur~$U=S_U$, et un $S_a$-isomorphisme $X_a\isomto
X_U\times_{S_U}S_a$. De plus, prenant~$U$ assez petit, on peut alors
supposer que l'isomorphisme~\eqref{eq:IX.4.*} provient d'un
isomorphisme:
$$
X_U'\isomto X_U\times_{S_U} S_U';
$$
ce dernier pourrait ne pas \^etre compatible avec les donn\'ees de
descente, cependant \`a condition de r\'etr\'ecir~$U$, il sera
compatible avec les donn\'ees de descente. Cela ach\`eve la
d\'emonstration.

\begin{corollaire}
\label{IX.4.5}
Sous les conditions de~\ref{IX.4.4}, pour que la donn\'ee de
descente sur~$X'$ soit effective, il faut et il suffit que pour tout
$a\in S$, la donn\'ee de descente correspondante sur~$X_a'$,
relativement au morphisme
$S_a'=S'\times_S\Spec(\cal{O}_a)\to\Spec(\cal{O}_a)$, le soit.
Lorsque~$S$ est localement noeth\'erien, et~$X'$ s\'epar\'e
sur~$S'$, on peut dans le crit\`ere pr\'ec\'edent remplacer
aussi~$\cal{O}_a$ par son compl\'et\'e.
\end{corollaire}

La
premi\`ere assertion r\'esulte de~\ref{IX.4.4} et~\ref{IX.4.3}, la
deuxi\`eme est alors cons\'equence de~\ref{IX.4.2}. Utilisant
encore~\ref{IX.4.2} et le fait que pour tout anneau local
noeth\'erien~$A$, on peut trouver un anneau local noeth\'erien
complet~$B$, et un homomorphisme local $A\to B$, tel que~$B$ soit plat
sur~$A$ et que $B/\goth{m}B$ soit une extension donn\'ee du corps
r\'esiduel $k=A/\goth{m}$ de~$A$, on trouve:

\begin{corollaire}
\label{IX.4.6}
Sous les conditions de~\ref{IX.4.4}, supposons de plus~$X'$
s\'epar\'e sur~$S'$, et~$S$ localement noeth\'erien. Pour que la
donn\'ee de descente sur~$X'$ soit effective, il faut et il suffit
que pour tout pr\'esch\'ema~$S_1$ sur~$S$, spectre d'un anneau
local complet \`a corps r\'esiduel alg\'ebriquement clos, la
donn\'ee de descente correspondante sur~$X_1'=X'\times_S S_1$,
relativement au morphisme $g_1\colon S_1'\to S_1$, soit effective.
\end{corollaire}

\begin{theoreme}
\label{IX.4.7}
Soit $g\colon S'\to S$ un morphisme fini et surjectif, et de
pr\'esentation finie (cette derni\`ere hypoth\`ese \'etant
cons\'equence des autres si~$S$ est localement
noeth\'erien)\footnote{On peut montrer qu'il suffit en fait que~$g$
soit un morphisme \emph{entier}, en se ramenant au cas du texte par un
proc\'ed\'e de passage \`a la limite dans le style EGA IV~8.}.
Alors~$g$ est un morphisme de descente effective pour la cat\'egorie
fibr\'ee des pr\'esch\'emas \'etales, s\'epar\'es, de type
fini sur d'autres.
\end{theoreme}

Il faut montrer que si~$X'$ est \'etale, s\'epar\'e, de type
fini sur~$S'$, et muni d'une donn\'ee de descente relativement
\`a~$g$, alors cette donn\'ee est effective.
Utilisant~\ref{IX.4.3}, on se ram\`ene facilement au cas o\`u~$S$
est noeth\'erien. Gr\^ace \`a~\ref{IX.4.5}, on peut donc
supposer que~$S$ est le spectre d'un anneau local noeth\'erien, a
fortiori que
$$
\dim S=n<+\infty.
$$
On raisonne alors par r\'ecurrence sur~$\dim S=n$, l'assertion
\'etant triviale pour $n<0$. Supposons donc $n\geq 0$ et le
th\'eor\`eme d\'emontr\'e pour les dimensions $n'<n$. En
vertu de~\ref{IX.4.6} on est ramen\'e au cas o\`u~$S$ est le
spectre d'un anneau local complet, donc~$S'$ est une r\'eunion finie
de spectres d'anneaux locaux complets. On a donc
$$
X'=X_1'\sqcup X_2'
$$
o\`u~$X_1'$ est \emph{fini} sur~$S'$, et o\`u~$X_2'$ n'a aucun
point au-dessus d'un des points ferm\'es
de~$S'$. Consid\'erons les morphismes
$$
\xymatrix@C=.5cm{{q_1,q_2\colon X''}\ar@<2pt>[r]\ar@<-2pt>[r]&{X'}}
$$
correspondants \`a la donn\'ee de descente, compatibles avec
$\xymatrix@C=.5cm{{p_1,p_2\colon S''}\ar@<2pt>[r]\ar@<-2pt>[r]&{S'}}$. On
voit aussit\^ot que
$$
X''=q_i^{-1}(X_1')\sqcup q_i^{-1}(X_2')\quad i=1,2
$$
est la d\'ecomposition canonique analogue de~$X''$ sur~$S''$, ce qui
implique $q_1^{-1}(X_1')= q_2^{-1}(X_1')$ et par suite~$X_1'$ et~$X_2'$
sont munis de donn\'ees de descente induites. Or soit~$T$ l'ouvert
de~$S$ compl\'ementaire de son point ferm\'e, donc $T'=S'\times_S
T$ est la partie de~$S'$ compl\'ementaire de l'ensemble des points
ferm\'es, et~$X_2'$, qui se trouve tout entier au-dessus de~$T'$,
est muni d'une donn\'ee de descente relativement au morphisme $T'\to
T$ induit par~$g$. Comme ce dernier est fini surjectif, et que $\dim
T<\dim S=n$, cette donn\'ee de descente est effective par
l'hypoth\`ese de r\'ecurrence. On voit donc qu'il suffit de
prouver que la donn\'ee de descente sur~$X_1'$ est effective, donc
on peut maintenant supposer~$X'$ \'etale et \emph{fini} sur~$S'$.
(N.B. le raisonnement par r\'ecurrence est inutile si on se borne
aux rev\^etements \'etales dans l'\'enonc\'e~\ref{IX.4.7}).
Soit alors~$S_0$ le spectre du corps r\'esiduel de~$A$, soit
$S_0'=S'\times_S S_0$ et d\'efinissons de m\^eme~$S_0''$, $S_0'''$
\`a partir des carr\'es et cubes fibr\'es~$S''$ et~$S'''$
de~$S'$ sur~$S$. En vertu de~\ref{IX.1.8}, les morphismes $S_0\to S$,
$S_0'\to S'$, etc. induisent des \'equivalences pour les
cat\'egories des rev\^etements \'etales de~$S$ et~$S_0$ d'une
part, $S'$ et~$S_0'$ d'autre part, etc. D'apr\`es les sorites de
la th\'eorie de la descente dans les cat\'egories \cite{IX.D}, il s'ensuit
que pour que $g\colon S'\to S$ soit un morphisme de descente effective
pour la cat\'egorie fibr\'ee des rev\^etements \'etales, il
faut et il suffit qu'il en soit ainsi de~$g_0\colon S_0'\to S_0$.
Mais c'est bien le cas, comme cas particulier de~\ref{IX.4.1}, par
exemple. Cela ach\`eve la d\'emonstration.

\begin{corollaire}
\label{IX.4.8}
La conclusion de~\ref{IX.4.7} subsiste si on suppose seulement que
$S'\to S$ est universellement submersif, de type fini et quasi-fini,
pourvu qu'on suppose~$S$ localement noeth\'erien.
\end{corollaire}

En
vertu de~\ref{IX.4.6}, on peut en effet supposer que~$S$ est le
spectre d'un anneau local noeth\'erien complet. Alors en vertu
de~\ref{IX.2.5}, il existe un morphisme fini et surjectif $S_1\to S$,
et un $S$-morphisme $S_1\to S'$. Comme $S_1\to S$ est un morphisme de
descente strict universel pour la cat\'egorie fibr\'ee
envisag\'ee, en vertu de~\ref{IX.4.7}, et que $S'\to S$ est un
morphisme de descente universel pour ladite, \ref{IX.4.8} r\'esulte
des sorites g\'en\'eraux~\cite{IX.D}.

\begin{corollaire}
\label{IX.4.9}
Soit $g\colon S'\to S$ un morphisme de type fini, surjectif et
universellement ouvert, avec~$S$ localement noeth\'erien. Alors~$g$
est un morphisme de descente effective pour la cat\'egorie
fibr\'ee des pr\'esch\'emas \'etales, s\'epar\'es et de
type fini sur d'autres.
\end{corollaire}

Proc\'edant comme dans~\ref{IX.4.7}, on est ramen\'e au cas
o\`u~$S$ est le spectre d'un anneau local noeth\'erien et
complet~$A$. Soit~$A_1$ une alg\`ebre finie sur~$A$, de
spectre~$S_1$, telle que $S_1\to S$ soit fini et \emph{surjectif},
donc un morphisme de descente effective universel pour la
cat\'egorie fibr\'ee envisag\'ee, gr\^ace \`a~\ref{IX.4.7}.
Il r\'esulte alors des th\'eor\`emes g\'en\'eraux \cite{IX.D}
que~$g$ est un morphisme de descente effective pour la cat\'egorie
fibr\'ee envisag\'ee, si et seulement si le morphisme
correspondant $g_1\colon S_1'=S'\times_S S_1\to S_1$ l'est. Comme ce
dernier satisfait aux m\^emes hypoth\`eses que~$g$, on est
ramen\'e \`a prouver~\ref{IX.4.9} pour~$S_1$ au lieu de~$S$.
Prenant d'abord pour~$A_1$ le compos\'e direct des $A/\goth{p}_i$,
pour les id\'eaux premiers minimaux~$\goth{p}_i$ de~$A$, on est
ramen\'e au cas o\`u~$A$ est \emph{intègre}. On montre
alors\footnote{\Cf EGA IV~14.3.13 et~14.5.4.} qu'il existe un
sous-sch\'ema int\`egre~$S_1$ de~$S'$, quasi-fini sur~$S$ et
dominant~$S$, passant par un point de la fibre de~$S'$ en le point
ferm\'e~$y$ de~$S$ (gr\^ace au fait que~$S'$ est universellement
ouvert de type fini sur~$S$ local noeth\'erien int\`egre, et
$S_y'\neq\emptyset$). Comme~$A$ est complet, $S_1$ est fini sur~$S$,
et comme il domine~$S$, le morphisme $S_1\to S$ est surjectif.
Rempla\c cant encore une fois~$S$ par~$S_1$, on est ramen\'e au
cas o\`u~$S'$ a une section sur~$S$, o\`u l'\'enonc\'e est
trivial.

\begin{theoreme}
\label{IX.4.10}
Soit $g\colon S'\to S$ un morphisme fini radiciel surjectif, de
pr\'esentation finie (cette derni\`ere condition \'etant
superflue si~$S$ est localement noeth\'erien\footnote{Il suffit
m\^eme que~$g$ soit entier, radiciel surjectif, comme on voit par
une r\'eduction facile au cas du texte, style EGA~IV~8, \cf SGA~4
VIII~1.1.}).
Alors le foncteur image inverse induit une \'equivalence de la
cat\'egorie des pr\'esch\'emas \'etales sur~$S$ avec la
cat\'egorie des pr\'esch\'emas \'etales sur~$S'$.
\end{theoreme}

Comme les morphismes diagonaux de~$S'$ dans $S'\times_S S'$ et
$S'\times_S S'\times_S S'$ sont des immersions surjectives, donc
induisent en vertu de~\ref{IX.1.9} des \'equivalences des
cat\'egories des pr\'esch\'emas \'etales sur~$S'\times_S S'$
\resp $S'\times_S S'\times_S S'$ avec la cat\'egorie des
pr\'esch\'emas \'etales sur~$S'$, il r\'esulte des sorites de
la descente \cite{IX.D} que tout~$X'$ \'etale sur~$S'$ est muni d'une
donn\'ee de descente et d'une seule relativement \`a $g\colon
S'\to S$. Donc~\ref{IX.3.3} implique que le foncteur image inverse
par~$g$, de la cat\'egorie des pr\'esch\'emas \'etales sur~$S$
dans la cat\'egorie des pr\'esch\'emas \'etales sur~$S'$, est
\emph{pleinement fidèle}. Reste \`a montrer qu'il est
essentiellement surjectif, \ie que tout~$X'$ \'etale sur~$S'$ est
isomorphe \`a l'image inverse d'un~$X$ \'etale sur~$S$. La
question \'etant \'evidemment locale sur~$S$ \emph{et sur}~$X'$,
on peut supposer $S$,~$S'$,~$X'$ affines. Mais alors~$X'$ est
s\'epar\'e de type fini sur~$S'$, et on peut appliquer le
crit\`ere d'effectivit\'e~\ref{IX.4.7}.

\begin{corollaire}
\label{IX.4.11}
La conclusion~\ref{IX.4.9} subsiste en rempla\c cant l'hypoth\`ese
sur~$g$ par: $g$~est fid\`element plat, quasi-compact et radiciel.
\end{corollaire}

M\^eme d\'emonstration, en invoquant~\ref{IX.4.1} au lieu
de~\ref{IX.4.7}.

On notera que la d\'emonstration de~\ref{IX.4.7} est
\og \'el\'ementaire\fg en ce qu'elle n'utilise pas les
th\'eor\`emes
de
finitude et de comparaison pour les morphismes propres (EGA~III~3,~4,~5). Il n'en est plus de m\^eme du r\'esultat suivant:

\begin{theoreme}
\label{IX.4.12}
Soit $g\colon S'\to S$ un morphisme propre, surjectif, de
pr\'esentation finie (cette derni\`ere hypoth\`ese \'etant
cons\'equence de la premi\`ere si~$S$ est localement
noeth\'erien). Alors~$g$ est un morphisme de descente effective
pour la cat\'egorie fibr\'ee des rev\^etements \'etales de
pr\'esch\'emas.
\end{theoreme}

En vertu de~\ref{IX.3.3} et de~\ref{IX.2.2}, on est ramen\'e \`a
prouver que pour tout rev\^etement \'etale~$X'$ sur~$S'$, muni
d'une donn\'ee de descente relativement \`a $g\colon S'\to S$,
cette donn\'ee de descente est effective. Utilisant~\ref{IX.4.3},
on est ramen\'e facilement
au cas o\`u~$S$ est noeth\'erien, et utilisant~\ref{IX.4.6}, on
peut donc supposer que~$S$ est le spectre d'un anneau local
noeth\'erien \emph{complet}~$A$. Introduisons~$S''$ et~$S'''$ comme
d'habitude, soit~$S_0$ le spectre du corps r\'esiduel de~$A$, et
soient~$S_0'$,~$S_0''$,~$S_0'''$ d\'eduits de~$S'$,~$S''$,~$S'''$
par le changement de base $S_0\to S$, \ie les fibres
de~$S'$,~$S''$,~$S'''$ au point ferm\'e de~$S$.
D'apr\`es~\ref{IX.1.10}, les morphismes $S_0\to S$, $S_0'\to S'$,
etc. induisent des \'equivalences de la cat\'egorie des
rev\^etements \'etales sur le sch\'ema-but avec la cat\'egorie
des rev\^etements \'etales sur le sch\'ema-source. Par suite,
$g\colon S'\to S$ est un morphisme de descente strict pour la
cat\'egorie fibr\'ee des rev\^etements \'etales de
pr\'esch\'emas, si et seulement si $g_0\colon S_0'\to S_0$ l'est,
ce qui est bien le cas en vertu de~\ref{IX.4.1}. Cela ach\`eve la
d\'emonstration de~\ref{IX.4.12}. (Dans ce raisonnement, on n'avait
besoin de~\ref{IX.1.10} que le fait que le foncteur envisag\'e
dans~\ref{IX.1.10} est \emph{pleinement fidèle}, ce qui n'utilise
\emph{pas} le th\'eor\`eme d'existence de faisceaux coh\'erents
en g\'eom\'etrie alg\'ebrique.)

\section{Traduction en termes du groupe fondamental}
\label{IX.5}

Soit
$$
g\colon S'\to S
$$
un \emph{morphisme de descente effective} pour la cat\'egorie
fibr\'ee des \emph{revêtements étales} de
pr\'e\-sch\'e\-mas, par exemple un morphisme propre, surjectif, de
pr\'esentation finie \eqref{IX.4.12}, ou un morphisme fid\`element
plat et quasi-compact. Introduisant comme d'habitude~$S''$,~$S'''$,
et d\'esignant par~$\cal{C}$,~$\cal{C}'$,~$\cal{C}''$,~$\cal{C}'''$
la cat\'egorie des rev\^etements \'etales
de~$S$,~$S'$,~$S''$,~$S'''$ respectivement, on a donc un diagramme
$2$-exact de cat\'egories
\begin{equation*}
\label{eq:IX.5.*}
\tag{$*$} 
\xymatrix@C=.5cm{{\cal{C}}\ar[r]^-{~p^*}&{\cal{C}'}\ar@<2pt>[rr]^{p_1^*,p_2^*}\ar@<-2pt>[rr]&&{\cal{C}''}\ar@<4pt>[rrr]^{p_{21}^*,p_{32}^*,p_{31}^*}\ar[rrr]\ar@<-4pt>[rrr]&&&{\cal{C}'''}}
\end{equation*}
correspondant au diagramme
$$
\xymatrix@C=.5cm{{S}&{S'}\ar[l]_-{~p}&&{S''}\ar@<2pt>[ll]\ar@<-2pt>[ll]_-{p_1,p_2}&&&{S'''.}\ar@<4pt>[lll]\ar[lll]\ar@<-4pt>[lll]_{p_{21},p_{32},p_{31}}}
$$
Supposons
les pr\'esch\'emas~$S$,~$S'$,~$S''$,~$S'''$ sommes disjointes de
pr\'esch\'emas connexes, ce qui sera le cas en particulier si ce
sont des pr\'esch\'emas localement connexes, a fortiori s'ils sont
localement noeth\'eriens (par exemple si~$S'$ est de type fini
sur~$S$ localement noeth\'erien). Alors les
cat\'egories~$\cal{C}$,~$\cal{C}'$ ... dans~\eqref{eq:IX.5.*} sont
des cat\'egories multigaloisiennes V~\ref{V.9}, d\'ecrites
chacune par une collection de groupes topologiques compacts totalement
disconnexes, savoir les groupes fondamentaux des composantes connexes
des pr\'esch\'emas~$S$,~$S'$,~$S''$,~$S'''$. Nous supposons pour
simplifier~$S$ connexe, et allons donner alors un proc\'ed\'e de
calcul pour son groupe fondamental, en termes de la cat\'egorie
fibr\'ee form\'ee avec~$\cal{C}'$,~$\cal{C}''$,~$\cal{C}'''$,
convenablement explicit\'ee \`a l'aide des groupes fondamentaux
exprimant ces cat\'egories. Le lecteur notera que le
proc\'ed\'e esquiss\'e est valable en fait dans le cadre
g\'en\'eral des cat\'egories multigaloisiennes (qui n'ont pas
\`a provenir de pr\'esch\'emas
donn\'es~$S$,~$S'$,~$S''$,~$S'''$). C'est d'ailleurs l'analogue du
proc\'ed\'e bien connu pour calculer le groupe fondamental d'un
espace topologique~$S$, r\'eunion localement finie de sous-espaces
ferm\'es~$S_i$ (ou r\'eunion quelconque de sous-espaces
ouverts~$S_i$), \`a l'aide des groupes fondamentaux des composantes
des~$S_i$
et des composantes des $S_i\cap S_j$. Bien entendu, la
situation analogue dans le cadre des pr\'esch\'emas tombe bien
dans le cadre g\'en\'eral de la descente, en introduisant le
pr\'esch\'ema~$S'$ somme des~$S_i$ et le morphisme canonique
$g\colon S'\to S$.

Posons
$$
E'=\pi_0(S'), \quad E''=\pi_0(S''), \quad E'''=\pi_0(S'''),
$$
o\`u~$\pi_0$ d\'esigne le foncteur \og ensemble des composantes
connexes\fg. Comme les produits fibr\'es de~$S'$ sur~$S$ forment un
objet simplicial de~$\Sch$, il est transform\'e par le
foncteur~$\pi_0$ en un ensemble simplicial dont~$E'$,~$E''$,~$E'''$
sont les composantes de dimension~$0$,~$1$,~$2$. Nous aurons \`a
utiliser les applications simpliciales
$$
q_i=\pi_0(p_i),\quad(i=1,2)\quad\text{et}
\quad q_{ij}=\pi_0(p_{ij}),\quad (i,j)=(2,1),(3,2),(3,1),
$$
mises en \'evidence dans le diagramme
\begin{equation*}
\label{eq:IX.5.1}
\tag{1} {\xymatrix@C=.5cm{{E'}&&{E''}\ar@<2pt>[ll]\ar@<-2pt>[ll]_{q_1,q_2}&&&{E'''.}\ar@<4pt>[lll]\ar[lll]\ar@<-4pt>[lll]_{q_{21},q_{32},q_{31}}}}
\end{equation*}
Les objets
de~$E'$ seront not\'es avec un accent, comme~$s'$, ceux
de~$E''$ \resp $E'''$ seront not\'es avec un~$''$ \resp un~$'''$.
Le fait que~$S$ soit connexe se traduit par $\pi_0(K)=0$, o\`u~$K$
est l'ensemble simplicial d\'efini par $g\colon S'\to S$, ou encore
par le fait que la relation d'\'equivalence dans~$E'$ engendr\'ee
par le couple d'applications $(q_1,q_2)$ est transitive.

Nous choisirons une fois pour toutes un \'el\'ement~$s_0'$
dans~$E'$, et pour tout~$s'$ dans~$E'$, un
\'el\'ement~$\overline{s'}\in E''$ tel que\footnote{On fera
attention que l'\'el\'ement~$\overline{s'}$ dont l'existence est
admise implicitement, n'existe pas dans tous les cas. Par suite, le
th\'eor\`eme~\ref{IX.5.1} tel qu'il est \'enonc\'e ne
s'applique pas dans tous les cas. Il n'est pas difficile cependant,
en s'inspirant du texte \'ecrit, de modifier l'\'enonc\'e de ce
th\'eor\`eme de telle fa\c con qu'il donne une m\'ethode de
calcul qui s'applique dans tous les cas. En particulier, les
corollaires dudit th\'eor\`eme sont valables tels quels.}
$$
q_1(\overline{s'})=s_0', \quad q_2(\overline{s'})=s',
$$
mettant ainsi en \'evidence la connexit\'e de~$S$. Pour
tout~$s'\in E'$, choisissons un point
g\'eom\'etrique~$\underline{s}'$ dans la composante connexe~$s'$
de~$S'$; ce point interviendra en fait par le foncteur-fibre~$F_{s'}'$
correspondant sur la cat\'egorie multigaloisienne~$\cal{C}'$. Le
groupe des automorphismes de ce foncteur, \ie le groupe fondamental
de~$S'$ en~$\underline{s}'$, sera not\'e~$\pi_{s'}$. On choisit de
m\^eme des~$\underline{s}''$ et des~$\underline{s}'''$, donc des
foncteurs~$F_{s''}''$ et~$F_{s'''}'''$, d'o\`u des groupes
fondamentaux~$\pi_{s''}$ et~$\pi_{s'''}$. Ainsi
$$
\pi_{s'}=\pi_1(S',\underline{s}'), \quad
\pi_{s''}=\pi_1(S'',\underline{s}''), \quad
\pi_{s'''}=\pi_1(S''',\underline{s}''').
$$
Pour tout $s''\in E''$, $p_1(\underline{s}'')$ se trouve dans la
m\^eme composante connexe que $\underline{q_1(s'')}$, donc il existe
un isomorphisme de foncteurs
$F''_{s''}\circ p_1^*\isomto F'_{s'}$
(\ie une \og classe de chemins\fg de~$p_1(\underline{s}'')$ \`a
$\underline{q_1(s'')}$). Cette remarque se r\'ep\`ete pour~$q_2$,
et les~$q_{ij}$. Choisissons toutes ces classes de chemins:
$$
F''_{s''}\circ p_i^*\isomto F'_{q_i(s'')},\quad F_{s'''}'''\circ
p_{ij}^*\isomto F'''_{q_{ij}(s''')},
$$
(pour
$i=1,2$ et $(i,j)=(2,1),(3,2),(3,1)$). Il en r\'esulte en
particulier des homomorphismes de groupes:
\begin{equation*}
\label{eq:IX.5.2}
\tag{2} {q_i^{s''}\colon \pi_{s''}\to\pi_{q_i(s'')},\quad
q_{ij}^{s'''}\colon \pi_{s'''}\to\pi_{q_{ij}(s''')},}
\end{equation*}
(m\^emes valeurs de~$i$ et $(i,j)$). Enfin, rappelons-nous que dans
la structure de cat\'egorie cliv\'ee de
fibres~$\cal{C}'$,~$\cal{C''}$,~$\cal{C'''}$ figurent aussi des
isomorphismes de foncteurs:
$$
p_{21}^*p_1^*\isomto p_{31}^*p_1^*,\quad p_{21}^*p_2^*\isomto
p_{32}^*p_1^*,\quad p_{31}^*p_2^*\isomto p_{32}^*p_2^*,
$$
d\'eduits d'isomorphismes des deux membres respectivement avec les
$u_i^*$~($i=1,2,3$), o\`u les~$u_i$ sont les trois projections
de~$S'''$ dans~$S'$. Quand on explicite ces donn\'ees, on trouve
pour tout~$s'''$ un \'el\'ement bien d\'etermin\'e
\begin{equation*}
\label{eq:IX.5.3}
\tag{3} {a_i^{s'''}\in\pi_{v_i(s''')},}
\end{equation*}
(o\`u les~$v_i$,~$i=1,2,3$ sont les trois applications $E'''\to E'$
d\'efinies par $v_i=\pi_0(u_i)$), d'ailleurs soumis aux conditions:
$$
q_1^{s_1''}q_{21}^{s'''}=\mathrm{int}(a_1^{s'''})
q_1^{s_2''}q_{31}^{s'''}\qquad (s_1''=q_{21}(s'''),\
s_2''=q_{31}(s''')),
$$
et les deux conditions analogues, faisant intervenir les~$a_2$
et~$a_3$. Le lecteur notera d'ailleurs que les
donn\'ees~\eqref{eq:IX.5.1},~\eqref{eq:IX.5.2},~\eqref{eq:IX.5.3}
permettent de reconstituer, \`a une \'equivalence de
cat\'egories fibr\'ees pr\`es, la cat\'egorie fibr\'ee
envisag\'ee de fibres~$\cal{C}'$,~$\cal{C}''$,~$\cal{C}'''$. Elles
doivent donc permettre en principe de reconstituer~$\cal{C}$ \`a
\'equivalence pr\`es, donc son groupe fondamental \`a
isomorphisme pr\`es. Nous d\'eterminerons en fait le groupe
fondamental en le point g\'eom\'etrique $p(\underline{s}_0')$
de~$S$, \ie le groupe des automorphismes de~$F'_{s_0'}\circ p^*$.

On note que la donn\'ee d'un objet~$X'$ de~$\cal{C}'$ est
\'equivalente essentiellement \`a la donn\'ee d'ensembles
finis~$X'_{s'}$ ($s'\in E'$) o\`u les~$\pi_{s'}$ op\`erent
contin\^ument. Une
donn\'ee de recollement sur un tel objet revient alors \`a la
donn\'ee, pour tout $s''\in E''$, d'une bijection:
$$
\varphi_{s''}\colon X'_{q_1(s'')}\isomto X'_{q_2(s'')}
$$
compatible avec les op\'erations de~$\pi_{s''}$, op\'erant sur
l'un et l'autre membre gr\^ace aux homomorphismes $q_i^{s''}\colon
\pi_{s''}\to\pi_{q_i(s'')}$. Prenant d'abord les~$s''$ de la
forme~$\overline{s'}$, on voit qu'une telle donn\'ee d\'efinit des
bijections
$$
\psi_{s'}\colon X'_{s_0'}=F_0'(X')\to X_{s'}'
$$
ce qui permet d'identifier les~$X'_{s'}$ au m\^eme ensemble
$F_0'(X')=X'_{s_0'}$, sur lesquels tous les groupes~$\pi_{s'}$ vont
d\`es lors op\'erer. Cela pos\'e, les
bijections~$\varphi_{s''}$ vont correspondre \`a des bijections
$$
g_{s''}\colon F_0'(X')\isomto F_0'(X'),
$$
soumis d'une part aux relations de commutation avec~$\pi_{s''}$:
\begin{equation*}
\label{eq:IX.5.a}
\tag*{a)} {g_{s''}q_1^{s''}(g'')=q_2^{s''}(g'')g_{s''}\qquad (s''\in E'',\ g''\in\pi_{s''}),}
\end{equation*}
d'autre part aux relations
\begin{equation*}
\label{eq:IX.5.b}
\tag*{b)} {g_{\overline{s'}}=g_{\overline{s_0'}}\qquad (s'\in E'),}
\end{equation*}
exprimant la fa\c con dont nous avions identifi\'e entre eux
les~$X_{s'}$. Quand on explicite la condition pour qu'une telle
donn\'ee de recollement soit en fait une donn\'ee de descente, on
trouve les relations:
\begin{equation*}
\label{eq:IX.5.c}
\tag*{c)} {a_3^{s'''}g_{q_{31}(s''')}a_1^{s'''}=
g_{q_{32}(s''')}a_2^{s'''}g_{q_{21}(s''')}\qquad (s'''\in E''').}
\end{equation*}

Cela nous donne une \'equivalence entre la cat\'egorie des objets
de~$\cal{C}'$ munis d'une
donn\'ee de descente, et la cat\'egorie des ensembles finis o\`u
les groupes~$\pi_{s'}$ op\`erent contin\^ument, munis de plus de
bijections~$g_{s''}$, satisfaisant les
relations~\ref{eq:IX.5.a},~\ref{eq:IX.5.b},~\ref{eq:IX.5.c}. Soit
alors~$G$ le groupe engendr\'e par les groupes~$\pi_{s'}$ et les
nouveaux g\'en\'erateurs~$g_{s''}$, soumis aux
relations~\ref{eq:IX.5.a},~\ref{eq:IX.5.b},~\ref{eq:IX.5.c}, et
soit~$\pi$ le groupe limite projective des quotients de~$G$ par les
sous-groupes d'indice fini dont les images inverses dans les
groupes~$\pi_{s'}$ soient des sous-groupes ouverts. On dit aussi
que~$\pi$ est le \emph{groupe de type galoisien engendré par
les~$\pi_{s'}$ et les~$g_{s''}$, soumis aux
relations~\ref{eq:IX.5.a},~\ref{eq:IX.5.b},~\ref{eq:IX.5.c}}. On
constate aussit\^ot que la cat\'egorie envisag\'ee est aussi
\'equivalente \`a la cat\'egorie des ensembles finis o\`u le
groupe topologique~$\pi$ op\`ere contin\^ument. Cela \'etablit
l'\'enonc\'e suivant:
\begin{theoreme}
\label{IX.5.1}
Soit $g\colon S'\to S$ un morphisme de pr\'esch\'emas qui soit un
morphisme de descente effective pour la cat\'egorie fibr\'ee de
rev\^etements \'etales de pr\'esch\'emas (\cf \ref{IX.4.9}
et~\ref{IX.4.12}). Supposons~$S$ connexe, et~$S'$, son carr\'e
fibr\'e~$S''$ et son cube fibr\'e~$S'''$, sommes de
pr\'esch\'emas connexes (ce qui est le cas par exemple si~$S'$ est
de type fini sur~$S$ localement noeth\'erien et connexe).
Choisissons comme dessus: un point g\'eom\'etrique dans toute
composante connexe de~$S'$,~$S''$,~$S'''$, certaines classes de
chemins, un $s_0'\in E'$, et pour tout $s'\in E'$ un
$s''\in E''$
dont les deux images dans~$E'$ soient~$s_0'$ et~$s'$.
($E'$,~$E''$,~$E'''$ d\'esignent respectivement l'ensemble des
composantes connexes de~$S'$,~$S''$,~$S'''$). Alors le groupe
fondamental de~$S$ en le point g\'eom\'etrique image de~$s_0'$ est
canoniquement isomorphe au groupe de type galoisien engendr\'e par
les $\pi_{s'}=\pi_1(S',\underline{s}')$ ($s'\in E'$) et des
g\'en\'erateurs~$g_{s''}$ ($s''\in E''$), soumis aux
relations~\ref{eq:IX.5.a},~\ref{eq:IX.5.b},~\ref{eq:IX.5.c} ci-dessus
faisant intervenir les \'el\'ements des groupes
$\pi_{s''}=\pi_1(S'',\underline{s}'')$, et les
\'el\'ements~$a_i^{s'''}$ ($i=1,2,3$, $s'''\in E'''$) introduits
plus haut.
\end{theoreme}

\begin{corollaire}
\label{IX.5.2}
Supposons que~$S'$ et~$S''$ n'aient qu'un nombre fini de composantes
connexes, et que les groupes fondamentaux des composantes connexes
de~$S'$ soient topologiquement de g\'en\'eration finie. Alors le
groupe fondamental de~$S$ est topologiquement de g\'en\'eration
finie.
\end{corollaire}

Ainsi,
nous prouverons plus tard que le groupe fondamental d'un sch\'ema
projectif normal sur un corps alg\'ebriquement clos est
topologiquement de g\'en\'eration finie. Utilisant le lemme de
Chow et la normalisation des sch\'emas alg\'ebriques, il
s'ensuivra que le m\^eme r\'esultat est vrai pour tout sch\'ema
propre sur un corps alg\'ebriquement clos.

\begin{corollaire}
\label{IX.5.3}
Supposons que~$S'$,~$S''$,~$S'''$ n'aient qu'un nombre fini de
composantes connexes, que les groupes fondamentaux des composantes
connexes de~$S'$ soient topologiquement de \emph{présentation
finie}, et les groupes fondamentaux des composantes connexes de~$S''$
topologiquement de \emph{génération finie}. Alors le groupe
fondamental de~$S$ est topologiquement de \emph{présentation
finie}.
\end{corollaire}

On notera qu'on peut exprimer~\ref{IX.4.9} (restreint aux
\emph{revêtements} \'etales) en disant qu'\emph{un morphisme
fini radiciel surjectif de préschémas noethériens induit
un isomorphisme des groupes fondamentaux}; de fa\c con imag\'ee,
on peut donc dire que le groupe fondamental est un \emph{invariant
topologique} pour les pr\'esch\'emas. On peut expliciter plus
g\'en\'eralement, \`a l'aide de~\ref{IX.5.1}, l'effet sur le
groupe fondamental d'op\'erations sur les pr\'esch\'emas, telles
que le \og pincement\fg du pr\'esch\'ema suivant un ensemble fini de
points, ayant une signification topologique simple. On trouve par
exemple:

\begin{corollaire}
\label{IX.5.4}
Soient $g\colon S'\to S$ un morphisme fini de pr\'esentation finie,
$T$ une partie discr\`ete de~$S$. Pour tout $s\in S$, soit $n(s)$ le
\og nombre g\'eom\'etrique de points\fg dans la fibre $g^{-1}(s)$
(qui s'explicite aussi comme le degr\'e s\'eparable de~$g^{-1}(s)$
sur~$\kres(s)$, somme des degr\'es s\'eparables de ses extensions
r\'esiduelles). On suppose que pour $s\in S-T$, on a $n(s)=1$. Pour
tout $s\in T$, soit $K_s$ une extension alg\'ebriquement close de
$\kres(s)$, $I_s$ l'ensemble des points g\'eom\'etriques de~$S'$ \`a
valeurs dans~$K_s$ (c'est un ensemble \`a $n(s)$ \'el\'ements),
$I_s'$ le compl\'ementaire d'un point choisi de~$I_s$, et enfin $I'$
l'ensemble r\'eunion des $I_s'$. On suppose $S'$ connexe. Alors le
groupe fondamental de~$S$ est isomorphe au groupe de type galoisien
engendr\'e par le groupe fondamental de~$S'$, et des
g\'en\'erateurs $g_i$ ($i\in I'$), soumis \`a aucune condition
suppl\'ementaire.
\end{corollaire}

Le
d\'etail de la d\'emonstration est laiss\'e au lecteur;
l'\'enonc\'e obtenu n'est que la traduction, en langage de la
th\'eorie des groupes, du fait qu'on a une \'equivalence de la
cat\'egorie $C$ des rev\^etements \'etales de~$S$, et de la
cat\'egorie des rev\^etements \'etales $X'$ de~$S'$,
\emph{munis} pour tout $s\in T$ d'un syst\`eme transitif de
bijections entre les $n(s)$ fibres de~$X'$ aux points de~$g^{-1}(s)$
\`a valeurs dans~$K_s$. (Sous cette forme intrins\`eque bien
entendu, il n'est plus n\'ecessaire de supposer $S'$ connexe).
\begin{exemple}
\label{IX.5.5}
On prouve facilement que la courbe rationnelle $\PP^1_k$ sur un corps
alg\'ebriquement clos $k$ est simplement
connexe\footnote{\Cf Exp.\, XI~\ref{XI.1.1}.}. Donc le groupe
fondamental d'une courbe rationnelle compl\`ete ayant exactement un
point double, \`a $n$ branches analytiques, est le groupe de type
galoisien libre engendr\'e par $n-1$ g\'en\'erateurs. Par
exemple, dans le cas d'un point double ordinaire, on trouve le groupe
fondamental $\hat\ZZ$, comme annonc\'e dans (I~\ref{I.11}~a)). Par
contre, l'existence d'un point de rebroussement (qui est un point
\og g\'eom\'etriquement unibranche\fg) n'a pas d'influence sur le
groupe fondamental.
\end{exemple}

\begin{corollaire}
\label{IX.5.6}
Soit $g\colon S'\to S$ un morphisme de pr\'esch\'emas
universellement submersif, \`a fibres g\'eom\'etriquement
connexes, $S$ \'etant connexe. Alors $S'$ est connexe, et
choisissant un point g\'eom\'etrique $s'$ dans~$S'$ et
d\'esignant par $s$ son image dans~$S$, l'homomorphisme
$$
\pi_1(S',s')\to \pi_1(S,s)
$$
est \emph{surjectif}. Si $g$ est un morphisme de descente effective
pour la cat\'egorie fibr\'ee des rev\^etements \'etales de
pr\'esch\'emas (\cf \emph{\ref{IX.4.12}}), introduisant le point
g\'eom\'etrique $s^{\prime\prime}=\diag(s')$
de~$S^{\prime\prime}=S'\times_S S'$, et les deux homomorphismes
$$
p_{1*},p_{2*}\colon \pi_1(S^{\prime\prime},s^{\prime\prime}) \to
\pi_1(S',s')
$$
induits par les deux projections, $\pi_1(S, s)$ est isomorphe au
conoyau de ce couple de morphismes dans la cat\'egorie des groupes
de type galoisien, \ie au quotient de~$\pi_1(S',s')$ par le
sous-groupe invariant ferm\'e engendr\'e par les
\'el\'ements de la forme
$p_{1*}(g^{\prime\prime})p_{2*}(g^{\prime\prime})^{-1}$, avec
$g^{\prime\prime}\in \pi_1(S^{\prime\prime},s^{\prime\prime})$.
\end{corollaire}

On sait en effet par~\ref{IX.3.4} que le foncteur $X\mto
X\times_SS'$ des rev\^etements \'etales sur~$S$ dans les
rev\^etements \'etales sur~$S'$ est pleinement fid\`ele, ce qui
\'equivaut au fait que l'homomorphisme sur les groupes fondamentaux
est un \'epimorphisme (V~\ref{V.6.9}). La derni\`ere assertion est
une cons\'equence imm\'ediate de la description~\ref{IX.5.1}.

\begin{remarque}
\label{IX.5.7}
Il n'est pas connu \`a l'heure actuelle si le groupe fondamental
d'un sch\'ema propre sur un corps alg\'ebriquement clos $k$ est
topologiquement de pr\'esentation finie\footnote{Cela semble
tr\`es improbable dans le cas des courbes lisses de genre $g\geq 2$,
en caract\'eristique $p>0$. Quand on remplace $\pi_1$ par son plus
grand quotient premier \`a $p$, par contre, il semble que les
techniques bien connues permettent de donner une r\'eponse
affirmative, m\^eme sans hypoth\`ese de propret\'e. \Cf un
travail en pr\'eparation de J.P.\, Murre.}. Utilisant~\ref{IX.5.3}, une
technique bien connue de sections hyperplanes, et la
d\'esingularisation des surfaces normales, on est ramen\'e au cas
d'une \emph{surface lisse sur} $k$. Cela permet du moins de montrer,
par voie transcendante, que la r\'eponse est affirmative en
caract\'eristique $0$ (et ceci sans \^etre oblig\'e d'admettre
la triangulabilit\'e de vari\'et\'es alg\'ebriques
singuli\`eres). En caract\'eristique $p>0$, la difficult\'e
principale semble dans le cas des courbes, dont on sait seulement que
le groupe fondamental est un quotient de celui qui se pr\'esente
dans le cas classique (\cf expos\'e suivant), le noyau par lequel on
divise \'etant cependant fort mal connu.
\end{remarque}

\begin{remarque}
\label{IX.5.8}
On pourrait expliciter d'autres cas particuliers que~\ref{IX.5.4}
et~\ref{IX.5.6} o\`u~\ref{IX.5.1} prend une forme
particuli\`erement simple. Un cas int\'eressant est celui o\`u
$S$ est le quotient de~$S'$ par un groupe fini d'automorphismes
$\Gamma$. \emph{Alors la catégorie des revêtements étales
de~$S'$ est équivalente à la catégorie des revêtements
étales $X'$ de~$S'$, où le groupe $\Gamma$ opère de fa\c con compatible avec ses opérations sur~$S'$, de telle fa\c con
que pour tout $s'\in S'$ et tout $g\in \Gamma_{s'}$} (o\`u
$\Gamma_{s'}$ d\'esigne \emph{le groupe d'inertie} de~$s'$
dans~$\Gamma$), \emph{$g$~opère trivialement dans la fibre}
$X_{s'}^\prime$. Si $S'$ est connexe cet \'enonc\'e
s'interpr\`ete de la fa\c con suivante. Soit $C'_0$ la
cat\'egorie des rev\^etements \'etales de~$S'$ o\`u $\Gamma$
op\`ere de fa\c con compatible avec ses op\'erations sur~$S'$
(mais sans satisfaire n\'ecessairement la condition ci-dessus sur
les groupes d'inertie des points de~$S'$). On voit facilement que
c'est une cat\'egorie galoisienne (V~\ref{V.5}), et que pour
tout point g\'eom\'etrique $a'$ de~$S'$, le foncteur fibre
$X'\mto X_{a'}'$ sur~$C'_0$ est un foncteur fondamental. Soit
$\pi_1(S',\Gamma;a')=G$ le groupe des automorphismes de ce foncteur,
muni de sa topologie habituelle. On a alors une suite exacte canonique
$$
e\to \pi_1(S',a')\to G\to \Gamma\to e
$$
(cas particulier de (V~\ref{V.6.13}), o\`u on prend pour $S$ le
rev\^etement trivial $S'\times \Gamma$ de~$S'$ d\'efini par
$\Gamma$, o\`u on fait op\'erer $\Gamma$ de fa\c con
\'evidente). D'ailleurs pour tout point g\'eom\'etrique $b'$
de~$S'$, on a un isomorphisme $\pi_1(S', \Gamma ; b' )\to G =
\pi_1(S', \Gamma; a')$ d\'efini \`a automorphisme int\'erieur
pr\`es provenant de~$\pi_1(S',a')$, et comme $\Gamma_{b'}$
s'applique de fa\c con \'evidente dans le premier membre, on
obtient un homomorphisme
$$
u_{b'}\colon \Gamma_{b'}\to G \quoi,
$$
d\'efini \`a automorphisme int\'erieur pr\`es (provenant
de~$\pi_1(S',a')$),
dont le compos\'e avec l'homomorphisme canonique
$G\to \Gamma$ est d'ailleurs l'immersion canonique $\Gamma_{b'} \to
\Gamma$. Ceci pos\'e, \emph{le groupe fondamental $\pi_1(S,a)$
\emph{est} canoniquement isomorphe au groupe quotient de~$G =
\pi_1(S', \Gamma; a')$ par le sous-groupe invariant fermé
engendré par les images des homomorphismes $\Gamma_{b'} \to
G$}. En particulier, l'image de~$\pi_1(S',a')$ dans $\pi_1(S,a)$ est
un sous-groupe invariant, et le quotient correspondant est isomorphe
\`a un quotient de~$\Gamma$. On peut d'ailleurs r\'eduire le
nombre des \og relations\fg introduites en introduisant, pour tout $g\in
\Gamma$, $ g\neq e$, le sous-pr\'esch\'ema $S'_g$ des
co\"incidences des automorphismes $\id_S$ et $g$ de~$S$, en
choisissant un point g\'eom\'etrique $b'_{g,i}$ dans chaque
composante connexe de~$S'_g$, puis un des homomorphismes
correspondants $\pi_1(S',\Gamma; b'_{g,i})\to G$, d'o\`u des
rel\`evements $\bar{g}_i$ de~$g$ dans~$G$. Il suffit alors de
prendre le quotient de~$G$ par le sous-groupe invariant ferm\'e
de~$G$ engendr\'e par les $\bar{g}_i$.

Lorsque $a'$ est invariant par $\Gamma$, on voit ais\'ement que
$\Gamma$ op\`ere de fa\c con naturelle sur~$\pi_1(S',a')$, et $G$
s'identifie au produit semi-direct correspondant. Identifiant alors~$\Gamma$
\`a
un sous-groupe de~$G$, on voit que dans les relations
introduites plus haut, faisant $b'=a'$, on trouve \og $g = e$\fg pour
$g\in I$. Donc \emph{si $S'$ a un point géométrique $a'$ fixe
par}~$I$ (\ie un point $s'$ dont le groupe d'inertie est $\Gamma$),
\emph{alors $\pi_1(S,a)$ est un groupe quotient du groupe quotient de
type galoisien de~$\pi_1(S',a')$ obtenu en \emph{\og rendant
triviales\fg} les opérations de~$\Gamma$ sur~$\pi_1(S',a')$; et il
est même isomorphe à ce dernier groupe si on suppose que pour
tout $g\in G$, l'ensemble d'inertie $S_g'$ est connexe, donc passe par
la localité de~$a'$}. Cette derni\`ere assertion est en effet
contenue dans la deuxi\`eme description donn\'ee plus haut pour
les relations \`a introduire dans $G$.

Ce dernier r\'esultat s'applique en particulier si l'on prend pour
$S'$ la puissance cart\'esienne $X^n$ d'un pr\'esch\'ema connexe
sur un corps alg\'ebriquement clos, pour $\Gamma$ le groupe
sym\'etrique $\Gamma={\goth {S}}_n$, op\'erant de la fa\c con
habituelle, d'o\`u pour $S$ la puissance sym\'etrique \nieme
de~$S$. Prenant alors pour $a'$ un point g\'eom\'etrique
localis\'e en la diagonale, on est sous les conditions
pr\'ec\'edentes, les ensembles d'inertie $S_g'$ contenant en effet
tous la diagonale. Utilisant le fait, prouv\'e dans l'expos\'e
suivant, que si $X$ est propre connexe sur~$k$, le groupe fondamental
de~$X^n$ s'identifie \`a $\pi_1(X)^n$, on trouve le r\'esultat
amusant suivant: \emph{Si $X$ est propre connexe sur~$k$
algébriquement clos, le groupe fondamental de sa puissance
symétrique \nieme, $n \geq 2$, est isomorphe au groupe
fondamental de~$X$ rendu abélien}. (J'ignore si le fait analogue
en Topologie alg\'ebrique est connu; il devrait pouvoir
s'\'etablir par la m\^eme m\'ethode de descente). Prenons par
exemple pour $X$ une courbe rationnelle $X=\PP^1_k$, on trouve une
\Nieme d\'emonstration du fait que $\PP^r_k$ est simplement
connexe, utilisant le fait que $\PP^1_k$ l'est. Prenons maintenant
pour $X$ une courbe simple sur~$k$, et $n\geq 2g-1$, de sorte que
$\Symm^n(X)$ est fibr\'e sur la jacobienne $J$, de fibres des
espaces projectifs, donc (comme on verra \`a l'aide des
r\'esultats des deux expos\'es suivants) a m\^eme groupe
fondamental que $J$. On retrouve alors sans d\'evissage le fait bien
connu que \emph{le groupe fondamental de la jacobienne de~$X$ est
isomorphe au groupe fondamental de~$X$ rendu abélien}.
\end{remarque}

\section[Une suite exacte fondamentale]{Une suite exacte fondamentale. Descente par morphismes \`a
fibres relativement connexes}
\label{IX.6}

\begin{theoreme}
\label{IX.6.1}
Soient $S$ le spectre d'un anneau artinien $A$ de corps r\'esiduel
$k$, $\overline{k}$ une cl\^oture alg\'ebrique de~$k$, $X$ un $S$
pr\'esch\'ema, $X_0=X\otimes_A k$, $\overline{X}_0=X\otimes_A
\overline{k}$, $\overline{a}$ un point g\'eom\'etrique
de~$\overline{X}$, $a$ son image dans~$X$, $b$ son image dans~$S$. On
suppose que $X_0$ est quasi-compact et \emph{géométriquement
connexe} sur~$k$ (N.B. si $X$ est propre sur~$S$, cela signifie que
$\H^0(X_0,\cal{O}_{X_0})$ est un anneau artinien \emph{local} de corps
r\'esiduel \emph{radiciel} sur~$k$). Alors la suite d'homomorphismes
canoniques
$$
e\mto \pi_1(\overline{X}_0,\overline{a})\to
\pi_1(X,a)\to\pi_1(S,b)\to e
$$
est exacte, et on a
$$
\pi_1(S,b)\isomfrom \pi_1(k,\overline{k})=\text{groupe de Galois
de~$\overline{k}$ sur~$k$.}
$$
\end{theoreme}

Comme les groupes fondamentaux ne changent pas en ch\^atrant par les
\'el\'ements nilpotents, on peut supposer $A=k$, ce qui rend
d\'ej\`a \'evident le dernier isomorphisme. Soit $k'$ la
cl\^oture s\'eparable de~$k$ dans~$\overline{k}$, et
consid\'erons $X'=X\otimes_k k'$, et l'image $a'$ de~$\overline{a}$
dans~$X'$. On a une suite d'homomorphismes canoniques
$$
e\mto \pi_1(\overline{X}_0,\overline{a})\to
\pi_1(X',a')\to\pi_1(S',b')\to e
$$
(o\`u $S'=\Spec(k')$). Enfin, on a un homomorphisme canonique de
cette suite dans celle relative \`a~$X/k$, gr\^ace au~diagramme
$$
\xymatrix{{S}&{X}\ar[l]&{\overline{X}_0}\ar[l]\\{S'}\ar[u]&{X'}\ar[u]\ar[l]&{\overline{X}_0  .}\ar[l]\ar@{=}[u]}
$$
On voit d'autre part que cet homomorphisme de suites de groupes est un
isomorphisme, comme il r\'esulte de~\ref{IX.4.11}. On est donc
ramen\'e \`a prouver que la deuxi\`eme suite est exacte, \ie on
peut supposer que $k$ est \emph{parfait}. Soient alors $k_i$ les
sous-extensions galoisiennes finies de~$k$ dans~$\overline{k}$, posons
$X_i=X\otimes_k k_i$, et soit $a_i$ l'image de~$\overline{a}$
dans~$X_i$. On laisse au lecteur
le soin
de v\'erifier que l'homomorphisme
naturel
$$
\pi_1(\overline{X}_0,a)\isomto \varprojlim\pi_1(X_i,a_i)
$$
est un isomorphisme, ce qui signifie simplement qu'un rev\^etement
\'etale de~$\overline{X}$ provient d'un rev\^etement \'etale
d'un~$X_i$, et que ce dernier est essentiellement unique, modulo
passage \`a un~$X_j$, $j\geq i$. D'autre part, soit $\pi_i$ le
groupe de Galois de~$k_i$ sur~$k$, \ie le groupe oppos\'e au groupe
des $S$-automorphismes de~$S_i=\Spec(k_i)$. Comme le foncteur
$S'\mto X\times_S S'$ des rev\^etements \'etales de~$S$ dans
les rev\^etements \'etales de~$X$ est pleinement fid\`ele
\eqref{IX.3.4}, il s'ensuit que $\pi_i$ est aussi isomorphe au groupe
oppos\'e aux groupes des $X$-automorphismes du rev\^etement
principal connexe $X_i$ de~$X$. Il r\'esulte donc de
V~\ref{V.6.13} que l'on a une suite exacte
$$
e\mto \pi_1(X_i,a_i)\to \pi_1(X,a)\to\pi_i\to e.
$$
Passant \`a la limite projective sur~$i$ dans ces suites exactes, on
trouve une suite exacte (puisqu'on est dans la cat\'egorie des
groupes de type galoisien), qui n'est autre que la suite envisag\'ee
dans~\ref{IX.6.1}. Cela ach\`eve la d\'emonstration.

La traduction de l'exactitude \`a droite dans~\ref{IX.6.1} en
langage g\'eom\'etrique est la suivante:

\begin{corollaire}
\label{IX.6.2}
Avec les notations pr\'ec\'edentes, soit $X'$ un rev\^etement
\'etale de~$X$, et soit $\overline{X'}_0$ le rev\^etement
\'etale correspondant de~$\overline{X}_0$. Les conditions suivantes
sont \'equivalentes:
\begin{enumerate}
\item[(i)] Il existe un $S'$ \'etale sur~$S$, et un $X$-isomorphisme
$X'\isomto X\times_S S'$, ($S'$ est alors d\'etermin\'e \`a
isomorphisme unique pr\`es en vertu de~\ref{IX.3.4}).
\item[(ii)] $\overline{X}_0'$ est compl\`etement d\'ecompos\'e
sur~$\overline{X}_0$.

Si $X'$ est connexe, ces conditions \'equivalent aussi \`a:
\item[(ii~bis)] $\overline{X}_0'$ a une section sur~$\overline{X}_0$.
\end{enumerate}
\end{corollaire}

(N.B. Ce dernier compl\'ement est essentiel; l'\'equivalence
de~(i) et (ii) signifie seulement que $\pi_1(S,b)$ est le groupe
quotient de~$\pi_1(X,a)$ par le sous-groupe
invariant ferm\'e engendr\'e par l'image
de~$\pi_1(\overline{X}_0,\overline{a})$, et non par cette image
elle-m\^eme). Sous les conditions pr\'ec\'edentes, nous dirons
que $X'$ est un rev\^etement \emph{géométriquement trivial}
\index{geometriquement trivial (revetement)@g\'eom\'etriquement trivial (rev\^etement)|hyperpage}de~$X$.

\begin{remarque}
\label{IX.6.3}
On ne peut dans l'\'enonc\'e~\ref{IX.6.1} remplacer $\overline{k}$
par une extension alg\'ebriquement close quelconque de~$k$, m\^eme
si $k$ est d\'ej\`a suppos\'e alg\'ebriquement clos. En
d'autres termes, il n'est pas vrai en g\'en\'eral que si $X$ est
un sch\'ema alg\'ebrique connexe sur un corps alg\'ebriquement
clos $k$, son groupe fondamental ne change pas en rempla\c cant $k$
par une extension alg\'ebriquement close; c'est d\'ej\`a faux
par exemple en caract\'eristique $p>0$ pour la droite affine
sur~$k$, \`a cause des ph\'enom\`enes de \og ramification
sup\'erieure\fg au point \`a l'infini, impliquant une structure
\og continue\fg pour le groupe fondamental. Nous verrons cependant dans
l'expos\'e suivant que de tels ph\'enom\`enes ne peuvent se
produire si $X$ est \emph{propre} sur~$k$. Nous montrerons aussi par
voie transcendante qu'il en est de m\^eme si $k$ est de
caract\'eristique nulle.
\end{remarque}

\begin{corollaire}
\label{IX.6.4}
Supposons que $a$ soit localis\'e en un $x\in X$ qui est rationnel
sur~$k$ (ou plus g\'en\'eralement, ayant un corps r\'esiduel
radiciel sur~$k$). Alors la suite exacte~\ref{IX.6.1} est scind\'ee.
\end{corollaire}

On peut supposer $S=\Spec(k)$. Si $x$ est rationnel sur~$k$, il
correspond \`a une section $S\to X$ de~$X$ sur~$S$, transformant $b$
en~$a$, et d\'efinissant un homomorphisme $\pi_1(S,b)\to \pi_1(X,a)$
qui est le scindage cherch\'e. Si $\kres(x)$ est radiciel sur~$k$, on se
ram\`ene au cas pr\'ec\'edent en faisant l'extension de la base
$\Spec(\kres(x))\to\Spec(k)$.

\begin{theoreme}
\label{IX.6.5}
Soient $f\colon X\to S$ un morphisme propre et surjectif de
pr\'esentation finie, \`a fibres g\'eom\'etriquement connexes,
$X'$ un pr\'esch\'ema de pr\'esentation finie et propre sur~$X$,
$s$ un point de~$S$, $F=X_s$ la fibre de~$X$ en~$s$, et $F_1'$ une
composante connexe de la fibre $F'=X'_s$ de~$X'$ en~$s$. Pour qu'il
existe un voisinage ouvert $X_1'$ de~$F_1'$ dans~$X'$, un
$S$-sch\'ema \'etale $S_1'$ et un $X$-isomorphisme $X_1'\isomto
S_1'\times_S X$, il faut et il suffit que $X'$ soit \'etale sur~$X$
aux points de~$F_1'$, et que $F_1'$ soit un rev\^etement
g\'eom\'etriquement trivial de~$F$.
\end{theoreme}

La
n\'ecessit\'e de la condition \'etant triviale, il reste
\`a prouver la suffisance. On se ram\`ene facilement au cas o\`u
$S$ est noeth\'erien. Consid\'erons la factorisation de Stein
$X\to T\to S$ de~$f$, o\`u $T$ est le spectre de l'Alg\`ebre
$f_*(\cal{O}_X)$ sur~$S$. Comme les fibres de~$X$ sur~$S$ sont
g\'eom\'etriquement connexes, et $f$ est surjectif, le morphisme
$T\to S$ est fini surjectif et radiciel, donc \eqref{IX.4.10} tout
$T'$ \'etale sur~$T$ provient par image inverse d'un $S'$ \'etale
sur~$S$. Cela nous ram\`ene \`a
prouver~\ref{IX.6.5}
en y rempla\c cant $S$ par~$T$, \ie dans le cas o\`u on suppose
$f_*(\cal{O}_X)=\cal{O}_S$. Consid\'erons alors la factorisation de
Stein $X'\to S'\to S$ du morphisme propre $h\colon X'\to S$, o\`u
$S'$ est le spectre de l'Alg\`ebre $h_*(\cal{O}_{X'})$. Les
morphismes $X'\to X$ et $X'\to S'$ d\'efinissent un morphisme
canonique
$$
X'\to X\times_S S',
$$
et notre assertion est contenue dans la suivante:

\begin{corollaire}
\label{IX.6.6}
Soit $f\colon X\to S$ un morphisme propre de pr\'esch\'emas
localement noeth\'eriens, tel que $f_*(\cal{O}_X)=\cal{O}_S$, et
soit $X'$ un pr\'esch\'ema propre sur~$X$. Consid\'erons la
factorisation de \emph{Stein} $X'\to S'\to S$ pour $X'\to S$, et le
morphisme canonique $X'\to X\times_S S'$. Soient $s$ un point de~$S$,
$s'$ un point de~$S'$ au-dessus de~$s$, correspondant \`a une
composante connexe $F_1'$ de la fibre $X'_s$ de~$X'$ en~$s$. Pour que
le morphisme $X'\to X\times_S S'$ soit un isomorphisme au-dessus d'un
voisinage ouvert $U'$ de~$s'$ \'etale sur~$S$, il faut et il suffit
que $X'$ soit \'etale sur~$X$ en les points de~$F_1'$, et que $F_1'$
soit un rev\^etement g\'eom\'etriquement trivial de la fibre
$F=X_s$.
\end{corollaire}

La n\'ecessit\'e \'etant encore triviale, il reste \`a prouver
la suffisance. La conclusion signifie aussi que a) le morphisme
d\'eduit de~$X'\to X\times_S S'$ par le changement de base
$\Spec(\widehat{\cal{O}}_{s'})\to S'$ est un isomorphisme, b) $S'$ est
\'etale sur~$S$ en~$s'$, \ie $\widehat{\cal{O}}_{s'}$ est \'etale
sur~$\widehat{\cal{O}}_s$. Sous cette forme, on voit que la conclusion
est invariante par changement de base $\Spec(\widehat{\cal{O}}_s)\to
S$. Les hypoth\`eses \'etant \'egalement stables par ce
changement de base, on peut donc supposer que $S$ est le spectre d'un
anneau local noeth\'erien complet. On peut de plus \'evidemment
supposer $X'$ connexe, ce qui implique ici que
$S'=\Spec(\cal{O}_{s'})$, $F'=F_1'$.

Comme
l'ensemble des points de~$X'$ o\`u $X'$ est \'etale sur~$X$ est
ouvert et contient la fibre $X'_{s'}=F'$, $X'$ \'etant propre
sur~$S$, il s'ensuit que $X'$ est \'etale sur~$X$. Comme il induit
sur~$F=X_s$ un rev\^etement \'etale isomorphe \`a un
$F\otimes_{\kres(s)}L$, o\`u $L$ est \'etale sur~$\kres(s)$, il
r\'esulte de~\ref{IX.1.10} qu'il est isomorphe \`a un
rev\^etement de la forme $X\times_S T$, avec $T$ \'etale sur~$S$.
(N.B. ici encore, il suffit d'utiliser que le foncteur
de~\ref{IX.1.10} est pleinement fid\`ele, qui r\'esulte du fait
qu'un isomorphisme formel de faisceaux coh\'erents sur~$X$ provient
d'un isomorphisme de ces faisceaux). Donc, si $T$ est d\'efini par
l'alg\`ebre~$B$ finie sur~$A$, $X'$ s'identifie au spectre de
l'Alg\`ebre $\cal{O}_X\otimes_A B$ sur~$X$, d'o\`u r\'esulte
aussit\^ot, puisque $f_*(\cal{O}_X)=\cal{O}_S$, que
$h_*(\cal{O}_{X'})$ est d\'efini par~$B$, donc l'homomorphisme
canonique $X'\to X\times_S S'$ n'est autre que l'isomorphisme
envisag\'e $X'\isomto X\times_S T$. Cela ach\`eve la
d\'emonstration.

\begin{corollaire}
\label{IX.6.7}
Sous les conditions de~\ref{IX.6.5}, pour qu'il existe un
pr\'esch\'ema $S'$ \'etale sur~$S$ et un $X$-isomorphisme
$X'\isomto X\times_S S'$, il faut et il suffit que $X'$ soit \'etale
sur~$X$ et que pour tout $s\in S$, $X'_{s'}$ soit un rev\^etement
g\'eom\'etriquement trivial de~$X_s$.
\end{corollaire}

En effet, s'il en est ainsi, $X'$ est r\'eunion d'ouverts $X_i'$ qui
sont isomorphes \`a des images inverses de~$S_i'$ \'etales
sur~$S$. On voit alors facilement que ces $S_i'$ se recollent en un
$S'$ \'etale sur~$S$, et qu'on obtient un isomorphisme $X'\isomto
X\times_S S'$. On peut dire par exemple que les $X_i'$ sont munis de
donn\'ees de descente relativement \`a~$X\to S$, qui se recollent
n\'ecessairement en une donn\'ee de descente sur~$X'$ tout entier
relativement \`a~$X\to S$. Et comme cette derni\`ere est effective
sur les~$X_i'$, il s'ensuit facilement (gr\^ace \`a un sorite
oubli\'e au num\'ero \ref{IX.4}) qu'elle est effective. On peut aussi
\'enoncer \ref{IX.6.7} sous la forme suivante:

\begin{corollaire}
\label{IX.6.8}
Soit $f\colon X\to S$ un morphisme propre surjectif de
pr\'esentation finie, \`a fibres g\'eom\'etriquement
connexes. Alors $f$ est un morphisme de descente effective pour la
cat\'egorie fibr\'ee des pr\'esch\'emas \'etales finis sur
d'autres. Le foncteur $S'\mto X\times_S S'$ induit une
\'equivalence de la cat\'egorie des pr\'esch\'emas \'etales
et finis sur~$S$ avec la cat\'egorie des pr\'esch\'emas
\'etales et finis sur~$X$ qui induisent sur chaque fibre $X_s$ un
rev\^etement g\'eom\'etriquement trivial.
\end{corollaire}

\begin{remarque}
\label{IX.6.9}
Soit $f\colon X\to S$ un morphisme propre et surjectif, avec $Y$
localement noeth\'erien, alors $f$ se factorise en un morphisme
$X\to S'$ satisfaisant l'hypoth\`ese de~\ref{IX.6.8}, et un
morphisme fini surjectif $S'\to S$ justiciable de~\ref{IX.4.7}, donc
en produit de deux morphismes qui sont des \emph{morphismes de
descente effective universels} pour la cat\'egorie fibr\'ee des
pr\'esch\'emas \'etales et finis sur d'autres. On peut en
conclure que $f$ lui-m\^eme est un morphisme de descente effective
universel pour la cat\'egorie fibr\'ee envisag\'ee. On retrouve
ainsi \ref{IX.4.12} par une m\'ethode diff\'erente.
\end{remarque}

\begin{remarque}
\label{IX.6.10}
La conclusion de~\ref{IX.6.7} ne reste pas valable si on remplace
l'hypoth\`ese que $f$ est propre par: $X$ est de type fini sur~$S$
et admet une section sur~$S$ (donc $f$ est universellement submersif
et un morphisme de descente pour la cat\'egorie fibr\'ee des
pr\'esch\'emas \'etales sur d'autres), m\^eme lorsque $S$ est
le spectre d'un anneau de valuation discr\`ete et lorsque $X'$ est
un rev\^etement \'etale de~$X$. Pour le voir, on part d'un $Z$
propre sur~$S$, dont la fibre g\'en\'erale
est une courbe rationnelle non singuli\`ere, et la fibre
sp\'eciale $Z_0$ consiste en deux droites concourantes. Par exemple,
si $t$ est une uniformisante de l'anneau de valuation~$A$, on prend le
sous-sch\'ema ferm\'e $Z$ de~$\PP_A^2$ d\'efini par
l'\'equation homog\`ene $x^2+y^2+tz^2=0$ (coordonn\'ees
homog\`enes $x$, $y$, $z$). On prend pour $X$ le compl\'ementaire
du point singulier $a$ de~$Z_0$ dans la r\'eunion $Z\cup
\PP_k^2$. Les fibres de~$X$ sont $\PP_k^1$ et $\PP_k^2 - a$, donc sont
g\'eom\'etriquement simplement connexes (\ie tout rev\^etement
\'etale d'une telle fibre est g\'eom\'etriquement trivial).
Cependant
on construit facilement, en proc\'edant comme dans le \No \ref{IX.4}, des
rev\^etements \'etales de~$X$ qui ne proviennent pas de
rev\^etements \'etales de~$S$, par recollement de rev\^etements
triviaux de~$Z-a$ et de~$\PP_k^2-a$. Il est possible par contre que la
conclusion de~\ref{IX.6.7} subsiste si on y remplace l'hypoth\`ese
de propret\'e par celle que $X$ soit \emph{universellement ouvert}
de pr\'esentation finie sur~$S$\footnote{C'est maintenant
prouv\'e, $g$ \'etant seulement universellement ouvert et
surjectif; \cf SGA~4 XV~1.15.}. C'est vrai du moins si on suppose que
les fibres de~$X$ sur~$S$ sont g\'eom\'etriquement
irr\'eductibles, et non seulement g\'eom\'etriquement connexes.
Signalons seulement qu'on peut dans cette question se ramener au cas
o\`u $S$ est le spectre d'un anneau de valuation discr\`ete
complet \`a corps r\'esiduel alg\'ebriquement clos.
\end{remarque}

L'interpr\'etation de~\ref{IX.6.7} en termes du groupe fondamental
est la suivante:

\begin{corollaire}
\label{IX.6.11}
Soit $f\colon X\to S$ un morphisme propre surjectif de
pr\'esentation finie, \`a fibres g\'eom\'etriquement
connexes. On suppose $X$ donc $S$ connexe. Soient $a$ un point
g\'eom\'etrique de~$X$, $b$ son image dans~$S$, et pour tout $s\in
S$, choisissons une cl\^oture alg\'ebrique $\overline{\kres(s)}$ de
$\kres(s)$, un point g\'eom\'etrique $a_s$ de~$S_s$ \`a valeurs
dans cette extension, et une classe de chemins de~$a_s$ dans~$a$,
d'o\`u un homomorphisme $\pi_1(\overline X_s,a_s)\to \pi_1(X,a)$,
o\`u $\overline X_s=X_s\otimes_{\kres(s)}\overline{\kres(s)}$. Alors
l'homomorphisme $\pi_1(X,a)\to \pi_1(S,b)$ est surjectif, et son noyau
est le sous-groupe invariant ferm\'e de~$\pi_1(X,a)$ engendr\'e
par les images des $\pi_1(\overline X_s,a_s)$.
\end{corollaire}

\begin{remarque}
\label{IX.6.12}
Sous les conditions de~\ref{IX.6.7}, supposant $S$ noeth\'erien,
on voit facilement que l'ensemble des points $s\in S$ tels que
$S'_{s'}$ soit g\'eom\'etriquement trivial sur~$X_s$ est un
ensemble constructible; si $X'$ est propre sur~$X$, il est m\^eme
ouvert, comme on voit sur~\ref{IX.6.6}. Cela permet donc, si $S$ et
un pr\'esch\'ema de Jacobson (par exemple de type fini sur un
corps), ou
$X'$ est propre sur~$X$, de se borner pour v\'erifier les conditions
de~\ref{IX.6.7} aux points $s$ de~$S$ qui sont \emph{fermés}. De
m\^eme, dans~\ref{IX.6.11} il suffit alors de prendre les
$\pi_1(\overline X_s,a_s)$ pour les points de~$S$ qui sont ferm\'es.
\end{remarque}

\begin{thebibliography}{D}{IX.7}

\bibitem[D]{IX.D} J.\, \textsc{Giraud}, M\'ethode de la
descente,
M\'emoire \no 2
de la Soc.\ Math.\ Fr.,~1964.

\bibitem[1]{IX.1} A.\, \textsc{Grothendieck}, G\'eom\'etrie Alg\'ebrique et
G\'eom\'etrie Formelle, S\'eminaire Bourbaki t.\, 11, 1959,
\No 182.
\end{thebibliography}

